% !PRINTED 407
\weberpage{407}
Jeder Teil der Summe auf der linken Seite der Formel (5),
der sich auf eine Idealklasse bezieht, gibt dann denselben Beitrag
zu der Summe, und wir erhalten:

\weq{7}{\begin{aligned}
\operatorname{Lim}_{s=1}(s-1)\sum\frac{1}{m^s}\sum\frac{(\Delta,n)}{n^s} &= h\,\frac{1}{\sqrt{\Delta}}\log\varepsilon, & \Delta&>0\\
&= h\,\frac{\pi}{\sqrt{-\Delta}}, & \Delta&<0.
\end{aligned}}

\gesperrt{Hier ist $h$ die Klassenzahl des Körpers $\Omega$, und um
nicht immer wieder die beiden Ausnahmefälle $\Delta = -4$,
$\Delta = -3$ anführen zu müssen, wollen wir festsetzen, daß
in diesen beiden Fällen, in denen die Klassenzahl 1 ist,
unter $h$ der Wert $\frac{1}{2}$ oder $\frac{1}{3}$ verstanden werden soll.}

So oft $n$ ein volles Restsystem nach dem Modul $\Delta$ durchläuft
, ist $\Sigma(\Delta,n) = 0$ [§ 85, (22)], und folglich bleibt die Summe
$\Sigma(\Delta,n)$, wenn die $n$ der Größe nach geordnet sind, wie weit
auch $n$ wächst, immer endlich. Danach ist der Satz Bd.\,II,
§ 196, 1. anwendbar, nach dem

\[
\operatorname{Lim}_{s=1}\sum\frac{(\Delta,n)}{n^s} = \sum\frac{(\Delta,n)}{n}
\]

ist. Andererseits ist

\[
\operatorname{Lim}(s-1)\Sigma\frac{1}{m^s} = 1
\]

(Bd.\,II, § 196, S. 723), und es ergibt sich aus (7):

\weq{8}{\begin{aligned}
h\log\varepsilon &= \sqrt{\Delta}\sum_{1,\infty}^{n}\frac{(\Delta,n)}{n}, & \Delta&>0\\
h\pi &= \sqrt{-\Delta}\sum_{1,\infty}^{n}\frac{(\Delta,n)}{n}, & \Delta&<0.
\end{aligned}}

Die Summe

\[
\sigma = \Sigma\frac{(\Delta,n)}{n}
\]

läßt sich in endlicher Form darstellen. Es ist nämlich nach
§ 86, (2):

\weq{9}{\sqrt{\Delta}\,(\Delta,n) = \overset{k}{\Sigma}(\Delta,k)\,e^{-\frac{2nk\pi i}{\Delta}},}

worin $k$ ein vollständiges Restsystem nach dem Modul $\Delta$ durchläuft
 und bei negativer Diskriminante $\sqrt{\Delta} = i\sqrt{-\Delta}$ zu setzen
ist. Danach folgt:

\[
\sigma = \frac{1}{\sqrt{\Delta}}\sum_{}^{k}(\Delta,k)\sum_{}^{n}\frac{1}{n}\,e^{-\frac{2nk\pi i}{\Delta}},
\]
%CONT>>
